%% ma: L12-B08-0155
%% lop: 12
%% chuong: 2
%% bai: 8
%% dang: 12.8.51
%% loai: TN4
%% muc: VD
%% dap_an: "C"
%% nguon: "CĐ2.4.pdf giáo viên gửi 27/10/2026 (M07, Lê Minh Khôi – Nguyễn Anh Khôi, Chương II Chủ đề 4 Các bài toán về tâm tỉ cự), Câu 5"
%% kien_thuc: "Bài 8 (biểu thức tọa độ của phép toán vectơ, độ dài vectơ), Bài 6 (điểm thỏa hệ thức vectơ, quy tắc ba điểm), Bài 7 (tọa độ vectơ, hình chiếu lên mặt phẳng tọa độ)"
%% trang_thai: chinh-thuc
%% ghi_chu: "Giáo viên đã duyệt 27/10/2026"
\begin{ex}
Trong không gian $Oxyz$, cho các điểm $A(-1;2;3)$, $B(6;-5;8)$ và $\overrightarrow{OM}=a\cdot\vec i+b\cdot\vec k$ với $a$, $b$ là các số thực thay đổi. Nếu $\left|\overrightarrow{MA}-2\overrightarrow{MB}\right|$ đạt giá trị nhỏ nhất thì giá trị của biểu thức $K=a-b$ bằng bao nhiêu?
\choice{$K=-25$}{$K=-13$}{\True $K=0$}{$K=26$}
\loigiai{$M(a;0;b)\in(Oxz)$. Gọi $I$ thỏa $\overrightarrow{IA}-2\overrightarrow{IB}=\vec 0\Rightarrow I=2B-A=(13;-12;13)$.

$\left|\overrightarrow{MA}-2\overrightarrow{MB}\right|=MI$ nhỏ nhất khi $M$ là hình chiếu của $I$ lên $(Oxz)$: $M(13;0;13)$. Vậy $K=0$.}
\end{ex}
